\section{One-step moment bounds}\label{sec:moments}

We now estimate the local map in environmental moments.  The distinction
between the environmental mean and fluctuations is essential: the mean
has square symmetry and gains the stronger pure transfer exponent,
whereas centered fluctuations gain a square-root saving from finite
dependence.  Rare large terms require the isolated integrations of
Section~\ref{sec:local-map}.

Fix the even moment order $P=40$.  For a stationary scale-$s$ field, put
\begin{equation}\label{mom:coordinates}
 \begin{split}
 X&=(\E\|f_0\|^P)^{1/P},\qquad
 Y=\|\E f_0\|,\qquad m=\E q_s(f_0),\\
 W&=\sum_{i\in\Z^2}
       \Cov\bigl(q_s(f_0),q_s(f_i)\bigr).
 \end{split}
\end{equation}
All norms are supremum norms over field variables; expectations in
\eqref{mom:coordinates} are environmental.  Finite dependence makes
the covariance sum finite.  Moreover $W\ge0$: if $\Lambda$ is a growing
square of cell indices, then
\[
 \frac1{|\Lambda|}\Var\!\left(\sum_{i\in\Lambda}q_s(f_i)\right)
 \longrightarrow W.
\]
Primes will denote the same coordinates for the output field at scale
$Ls$, in particular $q_{Ls}=Q/(Ls)$.

\begin{proposition}[One-step estimates]\label{mom:one-step}
One can choose $L$ large, then $b>0$ small, and finally $x_*>0$ such that
for every input field of Section~\ref{sec:local-map} with $X\le x_*$,
\begin{equation}\label{mom:bounds}
 \begin{split}
 |m'-Lm|&\le C_LX^2,\\
 Y'&\le\rho Y+C_L(|m|+X^2),\\
 X'&\le\rho X+C_L(\sqrt W+Y),
 \end{split}
\end{equation}
where $\rho<1$ and all constants are independent of the scale $s$.
The constants denoted by $C_L$ may depend on all the fixed step
parameters, including the subsequently chosen $b$.
\end{proposition}

We first establish the fluctuation estimate for the linearized map.
It is useful to isolate the dimension-free moment fact that permits
us to take a supremum over an exponentially large spin boundary.

\begin{lemma}[Moments of finitely dependent Hilbert vectors]
\label{mom:hilbert}
Let $(V_i)_{i\in I}$ be centered random vectors in a real Hilbert space,
where $|I|=N$.  Suppose a graph of bounded degree on $I$ is a dependency
graph: subfamilies with no graph edge between them are independent.
Write $a=\max_i(\E\|V_i\|^P)^{1/P}$ and $S=\sum_iV_i$.  Then
\begin{equation}\label{mom:hilbert-bound}
 \E\|S\|^P
 \le C_P(\E\|S\|^2)^{P/2}
       +C_{P,\mathrm{dep}}N^{P/2-1}a^P.
\end{equation}
The constants do not depend on the dimension of the Hilbert space.
In particular,
$\|S\|_{L^P}\le C_{P,\mathrm{dep}}N^{1/2}a$.
\end{lemma}

\begin{proof}
Write $P=2r$ and expand
$\|S\|^{2r}=(\sum_{i,j}\langle V_i,V_j\rangle)^r$.
For each tuple of $2r$ occurrences, connect two occurrences when their
indices coincide or are neighbors in the dependency graph.  A singleton
component has expectation zero by centering and independence.  If all
components have size at least two and one has size greater than two,
there are at most $r-1$ components.  Bounded graph degree gives at most
$C_{P,\mathrm{dep}}N^{r-1}$ such tuples: choose one index per component,
then the remaining indices along a spanning tree.  H\"older's inequality
bounds each expectation in absolute value by $a^{2r}$.

For the remaining tuples every component is a pair.  Independence
factors their expectation into contractions of pair covariance
operators.  Sum these contractions over all pairings and allow the
indices belonging to different pairs to be adjacent.  The added tuples
have merged components and incur the same
$C_{P,\mathrm{dep}}N^{r-1}a^{2r}$ error.  The unrestricted sum is exactly
the Wick expression for $\E\|Z\|^{2r}$, where $Z$ is a centered Gaussian
Hilbert vector with the covariance operator of $S$.  This operator is
positive even though individual cross-covariances need not be positive.
Diagonalizing it gives
$\E\|Z\|^{2r}\le C_P(\E\|Z\|^2)^r
=C_P(\E\|S\|^2)^r$.  Finite-dimensional projection and monotone
convergence give the same conclusion in a general Hilbert space.
Finally bounded dependence gives
$\E\|S\|^2\le C_{\mathrm{dep}}Na^2$.
\end{proof}

\subsection{Fluctuations of the linear map}

Fix a parent, translated to the origin, and transfer first to the
boundary of the square of radius $3Ls$.  Take the Hilbert norm there
with respect to the pure full-plane boundary marginal.  The pure
buffered density comparison, as stated in
\cite[Section~2, Proposition ``Buffered and small-patch comparisons'']{S},
bounds the final conditional transfer from this boundary to
$\partial I_0$ in supremum norm by a fixed constant times the Hilbert
norm.  Thus random Hilbert sums control the supremum in
\eqref{map:linear-formula}; the factor $p^{-1}A_0$ costs at most $p^{-1}$.
This factor is fixed independently of $L$.

Let
\[
 g_i=f_i-\E f_i,\qquad t_i=q_s(g_i).
\]
Choose the bounded energy representatives
$e_i=sP_{Cs,si}B_{si}$, enlarging the fixed support constant if necessary;
$B_{si}$ is the bond average of Theorem~\ref{pure:charge} at $si$,
and the transfer is through a square centered there.
Then $Q(e_i)=s$, $q_s(e_i)=1$, and $\|e_i\|\le C$.
Decompose
\[
 g_i=h_i+t_ie_i,\qquad Q(h_i)=0.
\]
Both $h_i$ and $t_i$ are centered in the environment, and
$\|h_i\|_{L^P(\E;\infty)}+\|t_i\|_{L^P(\E)}\le CX$.

The transferred $h_i$ have supremum norm at most
$CL^{-u}\|h_i\|$, by Theorem~\ref{pure:charge}.  One may first
transfer across a square centered at the child and then condition onto
the intermediate parent boundary.  There are $L^2$ children, with a
fixed dependence range in child units.  Lemma~\ref{mom:hilbert} bounds
their sum in $L^P(\E;L^2)$ by
\begin{equation}\label{mom:stable-fluctuation}
 CL^{1-u}X.
\end{equation}

For the charge part, let $w_i$ be the transfer of $e_i$ to the same
intermediate boundary.  These are deterministic Hilbert vectors with
\begin{equation}\label{mom:energy-vectors}
 \|w_i\|\le C/L,\qquad
 \|w_i-w_j\|\le CL^{-u}
 \quad\text{if }\|i-j\|_\infty\le2C.
\end{equation}
For the second estimate, $e_i-e_j$ has zero charge and is supported in
a fixed enlargement of one child neighborhood, so the zero-charge
transfer estimate applies.

Here the total covariance $W$ gives a sharper bound than the variance
of a single child.  Put $c(h)=\Cov(t_0,t_h)$ and let
$I=\mathcal C_0$.  Stationarity and finite dependence give the identity
\begin{align}
 \E\left\|\sum_{i\in I}t_iw_i\right\|^2
 &=W\sum_{i\in I}\|w_i\|^2\notag\\
 &\quad+\sum_{\substack{i\in I,h\\i+h\in I}}
 c(h)\langle w_i,w_{i+h}-w_i\rangle
 -\sum_{\substack{i\in I,h\\i+h\notin I}}c(h)\|w_i\|^2.
 \label{mom:covariance-regrouping}
\end{align}
Only a fixed number of $h$ contribute, and $|c(h)|\le CX^2$.
The first term is at most $CW$ because $W\ge0$ and
$\sum_i\|w_i\|^2\le C$.  By \eqref{mom:energy-vectors}, the second term
is $O(L^{1-u}X^2)$.  The last term has only $O(L)$ indices $i$ within
a fixed distance of the block boundary, so it is $O(L^{-1}X^2)$.
Consequently
\begin{equation}\label{mom:charge-variance}
 \E\left\|\sum_{i\in I}t_iw_i\right\|^2
 \le CW+C(L^{1-u}+L^{-1})X^2.
\end{equation}
Applying Lemma~\ref{mom:hilbert} with $N=L^2$ and $a=CX/L$ gives
\begin{equation}\label{mom:charge-P}
 \left\|\sum_{i\in I}t_iw_i\right\|_{L^P(\E;L^2)}
 \le C\sqrt W+
 C\bigl(L^{(1-u)/2}+L^{-1/2}+L^{-2/P}\bigr)X.
\end{equation}
The bound has no dependence on the dimension of the spin boundary.
Together with \eqref{mom:stable-fluctuation} and the final buffered
transfer, it implies
\begin{equation}\label{mom:linear-fluctuation}
 \|f'_{0,\mathrm{lin}}(f-\E f)\|_{L^P(\E;\infty)}
 \le\eta_LX+C\sqrt W,
 \qquad \eta_L\longrightarrow0.
\end{equation}
All constants in the limit statement use only the support range and
$p$, fixed before $L$.

\subsection{Charge, mean, and small Taylor remainders}

The exact normalization of $Q$ makes the first line of
\eqref{mom:bounds} especially simple.  Applying $q_{Ls}$ to
\eqref{map:linear-formula} integrates the new selection coins, so that
$p^{-1}$ cancels $\E_0A_j=p$.  Charge is preserved by pure transfer.
Therefore
\begin{equation}\label{mom:linear-charge}
 \E q_{Ls}(f'_{0,\mathrm{lin}})
 =\frac1L\sum_{i\in\mathcal C_0}\E q_s(f_i)=Lm.
\end{equation}

For the environmental mean, let
$\bar f_i=\E_{\mathrm{aux}}\E f_i$, where the auxiliary expectation
uses their independent reference law and leaves the spins fixed.
Then $\|\bar f_i\|\le Y$ and $q_s(\bar f_i)=m$.
The deterministic spin function $\bar f_i$ is quarter-turn invariant
about its center, because both the environmental distribution and the
auxiliary reference law respect this symmetry.  The transfer inside
$I_0$ integrates all these old auxiliary variables, so its mean is
the transfer of $\bar f_i$.  Subtract $m e_i$ from $\bar f_i$.
The difference has zero charge and norm at most $CY$,
since $|m|\le CY$.  Its transfer costs $CL^{-w}Y$ per child.  The charge
part costs $C_L|m|$.  Thus
\begin{equation}\label{mom:linear-mean}
 \|\E f'_{0,\mathrm{lin}}\|
 \le CL^{2-w}Y+C_L|m|.
\end{equation}
For use in the last line of \eqref{mom:bounds}, the simpler general
transfer bound also gives
\begin{equation}\label{mom:linear-total}
 \|f'_{0,\mathrm{lin}}\|_{L^P(\E;\infty)}
 \le\eta_LX+C_L(\sqrt W+Y).
\end{equation}

We record how flagged stencils affect low moments.  Enlarge a parent
stencil to include every comparison, threshold, and isolation check
that can affect its output, and put $S=\sum_{i\in\mathcal S}\|f_i\|$.
Its size depends only on the fixed step.  Let $\mathcal B$ be the event
that it contains a flag.  Markov's inequality and H\"older give
\begin{equation}\label{mom:bad-low-moments}
 \Prob(\mathcal B)\le C_{L,b}X^P,
 \qquad \E[S^r\1_{\mathcal B}]\le C_{L,b}X^P
 \quad(1\le r\le P).
\end{equation}
For example, for $r<P$ the second bound follows from
$(\E S^P)^{r/P}\Prob(\mathcal B)^{1-r/P}$; for $r=P$ it is immediate.
Choose $b$ after $L$ so small that $2|\mathcal S|b$ lies in the uniform
Taylor neighborhood of Lemma~\ref{map:linearization}.  If there is no
flag, every child norm is at most $2b$, and the regular Taylor remainder
$R$ satisfies
\begin{equation}\label{mom:regular-remainder}
 \|R\|\le C_LS^2\le C_LbS.
\end{equation}
In the second inequality the constant is independent of $b$.
The raw bound \eqref{map:raw-bound} and
\eqref{mom:bad-low-moments} show that the actual map differs from its
linearization by $O_L(X^2)$ in first environmental moment.  Equations
\eqref{mom:linear-charge} and \eqref{mom:linear-mean}, together with the
charge norm bound of Theorem~\ref{pure:charge}, now prove the first
two lines of \eqref{mom:bounds}, once $L$ is large enough that
$CL^{2-w}<1$.

The same reasoning does not yet prove a contraction in the $P$th
moment: the bound for $r=P$ in \eqref{mom:bad-low-moments} has size
$X^P$, not a higher power.  The next argument uses the location of the
flags to recover a small factor.

\subsection{The contribution of isolated large terms}

Consider one contributing child log $f_i$ in the raw output stencil
of \eqref{map:raw-bound}.  Enlarge the determining stencil to contain
the child-grid square of radius $100DL+2C$ about each such index $i$,
as well as every status check that can affect the output.  In particular,
this enlargement contains the full isolation neighborhood of any flag
group lying within $2C$ of a contributing log.  If a flag in this stencil
$k$ has $\|i-k\|_\infty>2C$, then $f_i$ is independent of that flag,
including its fresh threshold.  Hence
\begin{equation}\label{mom:independent-flag}
 \E[\|f_i\|^P\1_{\{k\text{ is a flag}\}}]
 =\E\|f_i\|^P\,\Prob(k\text{ is a flag})
 \le C_bX^{2P}.
\end{equation}
A union bound over the fixed enlarged stencil makes all such
contributions $O_{L,b}(X^2)$ in $P$th-root norm.

It remains to consider contributions for which every flag in this stencil
lies within $2C$ of $i$.  The flags then have diameter at most $4C$.
The inclusion of the full isolation neighborhoods ensures that they
form an isolated group of the construction.  All logs that can depend
on these flags lie within the $3d_0$ neighborhood transferred as $G_A$.
Equivalently, if isolation fails because of a distant additional flag,
that flag is independent of the contributing log and
\eqref{mom:independent-flag} applies.  This is the needed dichotomy;
small flag probability alone would not suffice.

For the isolated group, compare \eqref{map:isolated-ratio} with the
same expression after the other touching logs $B_A$ have been removed.
For every probability kernel $P_0$,
\begin{equation}\label{mom:remove-background}
 \left|\log\frac{P_0e^{B_A+G_A}}{P_0e^{B_A}}
                  -\log P_0e^{G_A}\right|
 \le2\|B_A\|.
\end{equation}
Each child log in $B_A$ lies outside the group neighborhood and is
independent of an actual triggering flag.  On the event that this group
occurs, its $P$th-root contribution is therefore $O_{L,b}(X^2)$ by
\eqref{mom:independent-flag}.  The same is true for every other log
omitted from the concentrated contribution.

After this removal, Lemma~\ref{pure:positive-weight} gives
\begin{equation}\label{mom:isolated-contraction}
 \left\|\log\E_0[e^{G_A}\mid I_A^c]
       -\E_0\log\E_0[e^{G_A}\mid I_A^c]\right\|
 \le \frac{C}{L}\sum_{i:\,\dist(i,A)\le3d_0}\|f_i\|.
\end{equation}
Indeed $G_A$ lies in a patch of radius a fixed multiple of $s$.
If this patch is not concentric with $I_A$, first transfer to the square
of radius $2Ls$ about the representative flag, which is contained in
$I_A$, and then condition onward.  Averaging a positive density cannot
increase its log oscillation.  Centering increases the oscillation
bound by at most a fixed factor.

There are a bounded number of possible flag patterns within a fixed
child-radius neighborhood of an anchor, and $O(L^2)$ possible anchoring
children for a given output stencil.  Separation bounds the number of
simultaneous isolated groups in that stencil by a constant independent
of $L$.  Raising \eqref{mom:isolated-contraction} to the $P$th power,
using the bounded number of logs in each group, and summing over anchors
therefore gives
\begin{equation}\label{mom:outlier-output}
 \|\1_{\mathcal B}f'_0\|_{L^P(\E;\infty)}
 \le CL^{2/P-1}X+C_{L,b}X^2.
\end{equation}
The first constant uses only the fixed child support radius; in
particular it is independent of $L$ and $b$.  This counting may be done
without disjointness of the anchoring events: use their union bound
before taking the $P$th root.

For completeness the linear output on the same event satisfies the
same bound, with a possibly different fixed constant.  Charge each
child term to a flag farther than $2C$ whenever one is present and use
\eqref{mom:independent-flag}.  Otherwise only a bounded child-radius
cluster can contribute a dependent large term.  The ordinary even
transfer to the parent boundary costs $C/L$ for that cluster.  The same
$O(L^2)$ anchor count then gives
\begin{equation}\label{mom:outlier-difference}
 \|\1_{\mathcal B}(f'_0-f'_{0,\mathrm{lin}})\|_{L^P(\E;\infty)}
 \le CL^{2/P-1}X+C_{L,b}X^2.
\end{equation}
Fixed multiplicities from the allocation of terms are included in the
constants.

\begin{proof}[Completion of Proposition~\ref{mom:one-step}]
The last line follows by combining \eqref{mom:linear-total},
\eqref{mom:regular-remainder}, and
\eqref{mom:outlier-difference}:
\[
 X'\le\bigl(\eta_L+CL^{2/P-1}+C_Lb\bigr)X
       +C_{L,b}X^2+C_L(\sqrt W+Y).
\]
Choose $L$ large enough that the first two coefficients, and the
coefficient $CL^{2-w}$ in \eqref{mom:linear-mean}, leave a strict margin
below one.  This is possible because $u>1$, $w>2$, and $P>2$.
Choose $b$ next to put flag-free stencils in the Taylor neighborhood
and make $C_Lb$ smaller than the remaining margin.  Finally choose
$x_*$ so that $C_{L,b}X$ is smaller than that margin for $X\le x_*$.
Increasing the resulting contraction factor to one common $\rho<1$
proves all three assertions.
\end{proof}

The mark parameters in Proposition~\ref{map:reversal} consequently
satisfy $\E p_j\le C_LX$.  We also record the precise Taylor consequence
needed below.

\begin{corollary}[Fourth-order moment substitution]
\label{mom:taylor-substitution}
Fix finitely many output stencils and a bounded multilinear expression
in at most four output potentials in total, allowing pure expectations,
environmental expectations, and products of such expectations.
The multilinear kernels are uniformly bounded in the scale; constants
may depend on these kernels and their fixed stencils.
Substitute the regular map's Taylor polynomial through degree four
and retain only total input degree at most four.  The resulting
environmental moment differs from that of the actual map by
$O_L(X^5)$, uniformly in $s$.
\end{corollary}

\begin{proof}
On the unflagged stencils, the fifth-order Taylor remainder and
H\"older's inequality give $O_L(X^5)$.  On flagged stencils use the
raw bound for the actual outputs and
\eqref{mom:bad-low-moments}.  Products of four fourth-order Taylor
polynomials have degree at most $16<P$, so the same bound controls
every term there.  Terms of input degree between five and sixteen
have expectation $O_L(X^5)$ when $X\le1$.  Products of separately
taken expectations obey the identical estimates.
\end{proof}
